\documentclass{article}
\usepackage{../../mystyle/Style}
\usepackage[margin=1in]{geometry}

\title{LaTeX Style Test}
\author{\defaultauthor}
\date{\today}

\begin{document}
\maketitle

\section{Theorem Environments}

\begin{definition}
    A group is a set $G$ equipped with a binary operation
    satisfying the group axioms.
\end{definition}

\begin{theorem}
    Every subgroup of a cyclic group is cyclic.
\end{theorem}

\begin{proof}
    Let $G = \langle g \rangle$ be a cyclic group.

    Suppose $H \leq G$.

    Then $H$ is cyclic.
\end{proof}

\begin{lemma}
    If $G$ is a group, then its identity element is unique.
\end{lemma}

\begin{proposition}
    Every cyclic group is abelian.
\end{proposition}

\begin{corollary}
    Every subgroup of a cyclic group is abelian.
\end{corollary}

\begin{example}
    The group $\Z_6$ is cyclic.
\end{example}

\begin{remark}
    All theorem environments should have upright body text.
\end{remark}


\section{Mathematical Commands}

\subsection{Number Sets}

\[
    \N \subseteq \Z \subseteq \Q \subseteq \R \subseteq \C
\]

\subsection{Absolute Values and Norms}

\[
    \abs{-5} = 5
\]

\[
    \norm*{\frac{x}{y}}
\]

\subsection{Sets}

\[
    A = \set{1,2,3}
\]

\[
    B = \set{x \in \R \mid x > 0}
\]

\subsection{Other Delimiters}

\[
    \ceil{3.2} = 4
\]

\[
    \floor{3.8} = 3
\]

\[
    \inner{x}{y}
\]

\subsection{Operators}

\[
    \ker(f), \quad
    \im(f), \quad
    \rank(A), \quad
    \Span\{v_1,v_2\}
\]

\subsection{Equations}

\begin{align*}
    (a+b)^2
    &= a^2 + 2ab + b^2 \\
    &= (a+b)(a+b)
\end{align*}

\begin{equation}
    \sum_{i=1}^{n} i = \frac{n(n+1)}{2}
\end{equation}


\section{Utility Commands}

\TODO{Review this theorem later.}

\[
    x \defeq 5
\]

This is \highlight{important information}.


\end{document}